\documentclass[11pt,openany]{book}
% ============================================================
%  PLANTILLA APUNTE-FORGE — adaptada de la plantilla minimalista
%  del usuario (article) a clase BOOK.
%
%  Decisión de diseño: cada TEMA generado por /generar-sintesis es
%  un \chapter. Con el tiempo, una misma materia acumula capítulos
%  y el documento se convierte literalmente en tu propio libro de
%  esa materia — no hace falta un documento nuevo por tema.
%
%  Mismo criterio estético que la plantilla original: gris oscuro
%  (no negro puro), títulos centrados en mayúsculas con línea
%  divisoria, márgenes estrechos, sin decoración innecesaria.
%  Los entornos nuevos (bloque, ejemplo, chequeo, ejercicio) usan
%  el mismo principio de minimalismo visual que rige todo el
%  contenido: una sola línea fina para separar, nunca cajas
%  pesadas ni colores de fondo — la separación es funcional
%  (distingue teoría / práctica / autoevaluación), no decorativa.
% ============================================================
\usepackage[utf8]{inputenc}
\usepackage[T1]{fontenc}
% NOTA: se remueve babel[spanish] por el mismo motivo que la
% plantilla original — el paquete de hifenación en español no está
% instalado en este entorno de compilación. Se renombran a mano los
% términos que babel normalmente traduciría (ver más abajo).
\usepackage{geometry}
\usepackage{amsmath}
\usepackage{amssymb}
\usepackage{amsthm}
\usepackage{enumitem}
\usepackage{tabularx}
\usepackage{booktabs}
\renewcommand{\familydefault}{\sfdefault}
\usepackage{titlesec}
\usepackage{fancyhdr}
\usepackage{ragged2e}
\usepackage[table]{xcolor}
\usepackage{array}
\usepackage{mdframed}

\geometry{
    top=2cm,
    bottom=2.2cm,
    left=2cm,
    right=2cm
}

% ---------- Colores (idénticos a la plantilla original) ----------
\definecolor{textgray}{RGB}{60, 60, 60}
\definecolor{rulegray}{RGB}{150, 150, 150}
\color{textgray}

% ---------- Metadatos editables por síntesis ----------
\newcommand{\miNombre}{Nombre Apellido}      % quién estudia (footer)
\newcommand{\materiaActual}{Materia}          % se resetea por materia
\newcommand{\setMateria}[1]{\renewcommand{\materiaActual}{#1}}

% ---------- Header / Footer (igual que el original) ----------
\pagestyle{fancy}
\fancyhf{}
\renewcommand{\headrulewidth}{0pt}
\renewcommand{\footrulewidth}{0pt}
\fancyfoot[C]{\small\color{textgray}\thepage}
\fancyfoot[R]{\footnotesize\color{textgray}\miNombre}
\fancyfoot[L]{\footnotesize\color{textgray}\materiaActual}

% ---------- Traducciones manuales (sin babel) ----------
\renewcommand{\tablename}{Tabla}
\renewcommand{\figurename}{Figura}
\renewcommand{\contentsname}{Índice}
\newcommand{\ok}{\checkmark}

% ============================================================
%  CAPÍTULO = TEMA. Mismo criterio visual que \section en la
%  plantilla original (centrado, mayúsculas, línea fina), pero un
%  punto más grande porque es el nivel jerárquico superior acá.
% ============================================================
\titleformat{\chapter}
  {\centering\color{textgray}\fontsize{18}{21}\selectfont\bfseries}
  {}
  {0pt}
  {}
  [\vspace{3pt}\color{rulegray}\titlerule\vspace{6pt}]
\titlespacing*{\chapter}{0pt}{0pt}{18pt}
% NOTA: a diferencia de \section (donde redefinir el comando funciona sin
% problema), redefinir \chapter directamente rompe el mecanismo interno de
% \tableofcontents en la clase book (corrompe la entrada de TOC). En vez de
% eso, se usa un comando auxiliar explícito: \temacapitulo{<Tema>} en lugar
% de \chapter{<Tema>} directo.
\newcommand{\temacapitulo}[1]{\chapter{\MakeUppercase{#1}}}

% ============================================================
%  SECTION = cada uno de los 8 pasos de la plantilla pedagógica
%  (1. Motivación, 2. Mapa de fuentes, ... 8. Ejercicio guiado).
%  Formato idéntico al de la plantilla original.
% ============================================================
\titleformat{\section}
  {\centering\color{textgray}\fontsize{13}{15}\selectfont\bfseries}
  {}
  {0pt}
  {}
  [\vspace{2pt}\color{rulegray}\titlerule\vspace{4pt}]
\titlespacing*{\section}{0pt}{18pt}{6pt}
\let\oldsection\section
\renewcommand{\section}[1]{\oldsection{\MakeUppercase{#1}}}

% ============================================================
%  SUBSECTION = micro-bloques dentro de "3. Marco teórico"
%  (3.a definición, 3.b prompt generativo — 3.c usa el entorno
%  \chequeo, no subsection, porque necesita numeración propia).
% ============================================================
\titleformat{\subsection}
  {\color{textgray}\normalsize\bfseries}
  {}
  {0pt}
  {}
\titlespacing*{\subsection}{0pt}{12pt}{4pt}

% ============================================================
%  ENTORNOS CUSTOM — uno por cada pieza funcional de la
%  plantilla pedagógica. Numerados por capítulo (=por tema).
% ============================================================

\theoremstyle{definition}
\newtheorem{bloqueteorico}{Bloque}[chapter]
% Paso 3.a — Definición/fórmula del bloque, con notación de la fuente.
% Uso: \begin{bloque}[Fuente: Lin \& Costello]{Distancia de Hamming} ... \end{bloque}
\newenvironment{bloque}[2][]{%
  \bloqueteorico%
  \noindent\textbf{\color{textgray}#2}%
  \ifx&#1&\else\ {\footnotesize\color{rulegray}(#1)}\fi%
  \par\vspace{2pt}\color{rulegray}\hrule\color{textgray}\vspace{4pt}%
}{%
  \vspace{4pt}
}

\theoremstyle{plain}
\newtheorem{ejemploresueltoct}{Ejemplo resuelto}[chapter]
% Paso 4 — Ejemplo resuelto. Argumento opcional: "completo" | "parcial"
% (alternar entre ambos formatos por regla de pedagogia-cognitiva).
\newenvironment{ejemplo}[1][completo]{%
  \ejemploresueltoct%
  \ifthenelse{\equal{#1}{parcial}}%
    {\ {\footnotesize\itshape\color{rulegray}(parcial — completá el paso final)}}%
    {}%
  \par\vspace{2pt}\color{rulegray}\hrule\color{textgray}\vspace{4pt}%
}{%
  \vspace{4pt}
}

% Paso 3.c — Chequeo de comprensión (adjunct questions). Va INMEDIATAMENTE
% después de cada bloque teórico, nunca concentrado al final del capítulo.
\newmdenv[
  linecolor=rulegray,
  linewidth=0.4pt,
  topline=true,
  bottomline=true,
  leftline=false,
  rightline=false,
  innertopmargin=6pt,
  innerbottommargin=6pt,
  skipabove=8pt,
  skipbelow=8pt
]{chequeobox}
\newenvironment{chequeo}{%
  \begin{chequeobox}%
  {\footnotesize\bfseries\color{rulegray} CHEQUEO RÁPIDO}\par\vspace{2pt}%
  \begin{enumerate}[leftmargin=1.4em, itemsep=2pt, label=\arabic*.]
}{%
  \end{enumerate}
  \end{chequeobox}
}

% Paso 6.a / 6.b — Ejercicios. Argumentos: tipo (bloque|intercalado) y
% familia, para que quede trazable contra mapa-estudio.json a simple vista.
\newtheorem{ejerciciocount}{Ejercicio}[chapter]
\newenvironment{ejercicio}[2][]{%
  \ejerciciocount%
  \ifx&#2&\else\ {\footnotesize\color{rulegray}[#2]}\fi%
  \ifx&#1&\else\ {\footnotesize\itshape\color{rulegray}— familia: #1}\fi%
  \par\vspace{2pt}
}{%
  \vspace{6pt}
}

% Paso 5 — Errores comunes. Lista simple, sin caja (el tono neutro va
% en el TEXTO, no en la decoración — coherente con minimalismo visual).
\newenvironment{erroresComunes}{%
  \begin{itemize}[leftmargin=1.4em, itemsep=4pt, label={\color{rulegray}\textbullet}]
}{%
  \end{itemize}
}

% Paso 1 — Motivación. Estilo narrativo en primera persona, apertura
% con oración temática (itálica, con leve sangría) antes del cuerpo.
\newenvironment{motivacion}{%
  \begin{quote}
  \itshape\color{textgray}
}{%
  \end{quote}
}

% Necesario para el \ifthenelse de \ejemplo
\usepackage{ifthen}



\setMateria{Comunicaciones Digitales}

\begin{document}

\begin{titlepage}
  \centering
  \vspace*{4cm}
  {\fontsize{22}{26}\selectfont\bfseries\color{textgray} \materiaActual \par}
  \vspace{0.5cm}
  {\large\color{rulegray}\rule{0.85\linewidth}{0.4pt}\par}
  \vspace{1cm}
  {\large\color{textgray} Síntesis teórica personal \par}
\end{titlepage}

\tableofcontents

\temacapitulo{Decodificadores convolucionales}

\section{Motivación}
\begin{motivacion}
Estás diseñando un enlace de comunicación y el canal tiene ruido — algunos bits van a llegar mal sí o sí. La pregunta no es cómo evitar el error, es cómo diseñar un código que te permita reconstruir el mensaje original aunque una parte de él llegue corrompida.
\end{motivacion}

Este es exactamente el problema que resuelven los códigos convolucionales: agregar redundancia estructurada para poder decodificar incluso con errores de transmisión.

\section{Mapa de fuentes}
Lin \& Costello desarrollan el formalismo algebraico completo del código convolucional. El paper original de Viterbi (1967) aporta el algoritmo de decodificación específico. Ambas fuentes usan notación compatible, sin contradicciones de fondo detectadas para este tema.

\section{Marco teórico}

\begin{bloque}[Fuente: Lin \& Costello]{Polinomio generador}
Un código convolucional se define mediante polinomios generadores $g_1(x), g_2(x), \ldots$ que determinan cómo se combinan los bits de entrada con la memoria del codificador para producir la salida codificada.
\end{bloque}

\begin{chequeo}
\item ¿Qué pasa con la tasa de código si agregás un polinomio generador más?
\item ¿Por qué la memoria del codificador afecta la distancia libre del código?
\end{chequeo}

\begin{bloque}[Fuente: Viterbi, 1967]{Algoritmo de Viterbi}
Un algoritmo de programación dinámica que encuentra el camino de máxima verosimilitud a través del diagrama de trellis del código, sin necesidad de evaluar todos los caminos posibles explícitamente.
\end{bloque}

\begin{chequeo}
\item ¿Por qué el algoritmo de Viterbi es más eficiente que fuerza bruta sobre todos los caminos del trellis?
\end{chequeo}

\section{Ejemplo resuelto}

\begin{ejemplo}[completo]
Dado un codificador con polinomios generadores $g_1 = 111$, $g_2 = 101$, y una secuencia de entrada $u = 1011$, calculamos la salida codificada paso a paso combinando cada bit de entrada con la memoria actual del registro de desplazamiento.
\end{ejemplo}

\begin{ejemplo}[parcial]
Con el mismo codificador, la entrada $u = 1101$ produce una salida codificada donde los primeros 3 pares de bits son $(1,1), (0,1), (1,0)$ — completá el último par aplicando el mismo procedimiento.
\end{ejemplo}

\section{Errores comunes}
\begin{erroresComunes}
  \item Confundir la tasa de código con la distancia libre — son propiedades distintas del mismo código.
  \item Asumir que más memoria siempre implica mejor corrección de errores sin considerar el costo computacional del decodificador.
\end{erroresComunes}

\section{Ejercicios}

\begin{ejercicio}[decodificación de secuencias]{6.a — bloque}
Calculá la salida codificada para el codificador $g_1=111, g_2=101$ con entrada $u=0110$.
\end{ejercicio}

\begin{ejercicio}[decodificación de secuencias]{6.b — intercalado}
Mezclando con distancia de Hamming (tema previo): dadas dos secuencias codificadas candidatas, determiná cuál está más cerca de la recibida usando distancia de Hamming, y after explicá por qué el algoritmo de Viterbi hace ese mismo cálculo de forma más eficiente.
\end{ejercicio}

\end{document}
